%
% $Log: sentences.tex,v $
% Revision 1.4  2002/10/21 15:26:44  mjm
% to ohio, revised
%
% Revision 1.3  2001/08/15 18:07:03  mjm
% APPM preprint 466
%
% Revision 1.2  1999/08/09 19:55:27  mjm
% beta release
%
% Revision 1.1  1999/06/10 23:34:48  mjm
% Initial revision
%
%

\handout{Flow}
       
Written mathematics must be readable. This may seem trivial, but it is
an important point. You should be able to read your work aloud to a
classmate and have them understand your solution. If you need to add
any explanations, these should be included in your written work.
The most common mistake is to write mathematics without using enough
words. All writing, even mathematics, should consist of complete sentences.
These should explain the problem by providing both the method and 
justification for each  step of the solution.


\begin{flushleft}
{\bf{Why are sentences important in mathematics?}}
\end{flushleft}

  Although sometimes it seems hard to read textbooks, it would be
much harder to understand if they only had equations and no sentences.
The situation is similar in lecture: if the professor just listed
formulas on the chalkboard without talking about them, leaving you to
figure out what was being done in each step, how much could you
understand from the lecture?  Neither of these would be a good way for
most students to learn, since sentences are necessary to explain the
mathematics.

\begin{flushleft}
  \bf{Why should students use sentences in a Mathematics class?}
\end{flushleft}

In a Mathematics class, you should explain homework solutions using complete
sentences.  That means linking together thoughts with words and
embedding equations into sentences.  Going through the extra work to
do this will benefit you in several ways:

\begin{itemize}
\item Writing down your thoughts and organizing them into complete
  sentences will help you to understand the method of solution better.
  
\item When you look back on homework to study for a test, or later on
  in another class, you will understand what you
  were doing on each problem and the mathematics behind it.
  
\item Other people (teacher, classmates, grader,...) will understand
  what you are doing at each step, and why you are doing it.  This
  way, you won't lose points for skipping steps or solving the problem
  in an unusual way.
  
\item Communicating your work will be essential in whatever field you choose.
  Even though the fields are stereotypically weak on
  writing, engineers and scientists spend a surprising amount of time
  writing reports and giving oral presentations.
\end{itemize}

%%%%%%%%%%%%%%%%%%%%%%%%%%%%%%%%%%%%%%%%%%%%%%%%%%%%%%%%%%%%%%%%%%%%%%%%%
\subsubsection*{The Royal ``We''.}

It is customary to write mathematics using ``we'' instead of ``I'' and
``you''. Think of yourself as the tour guide, showing the sights to
your reader. You and the reader together make ``we''.


\subsubsection*{Your Audience}
Keep in mind who your audience is. For the Good Problems, pretend that
you are writing for a classmate who has not seen this problem
before. Write with enough detail so that they can follow your
explanations.  A good way to test if you have written enough is to read
your work aloud to another person. If you need to add any words to
make it make sense, then these words should be included in your
written work.



%%%%%%%%%%%%%%%%%%%%%%%%%%%%%%%%%%%%%%%%%%%%%%%%%%%%%%%%%%%%%%%%%%%%%%%%%
\secondpage
\subsection*{{Examples:}}

\begin{itemize}

%\noindent{\bf How to put an equation into a sentence:}
%\subsubsection*{How to put an equation into a sentence:}
%
%\begin{goodexample}
%    The derivative of the curve $y= 3\sqrt{x}$ is $\frac{dy}{dx} =
%    \frac{3}{2 \sqrt{x}}$.
%\end{goodexample}  
%    
%\begin{badexample}
%  The derivative $\frac{dy}{dx}$ is $\frac{3}{2 \sqrt{x}}$ for the
%  curve $y$ which is $3 \sqrt{x}$.
%\end{badexample}
%    
%\begin{badexample}
%  $\frac{dy}{dx}$ is the derivative of $3 \sqrt{x}$ which is
%  $\frac{3}{2 \sqrt{x}}$.
%\end{badexample}



%\bigskip
%%%%%%%%%%%%%%%%%%%%%%%%%%%%%%%%%%%%%%%%%%%%%%%%%%%%%%%%%%%%%%%%%%%%%%%%%
%\subsubsection*{Using sentences incorrectly:}

\item  
Find an equation for the tangent to the curve $y = x + \frac{2}{x} $
at the point $(1,3)$.
 
%\begin{badexample}
%
%  \begin{displaymath}
%    \frac{dy}{dx} = 1 - \frac{2}{x^2}
%  \end{displaymath}
%                                %
%  so the equation for the derivative is
  %
%  \begin{displaymath}
%    \frac{dy}{dx}= -1 = m
%  \end{displaymath}
%                                %
%  so the slope at $(x,y)$ can be found by $y-3 = -1 (x-1)$ so $y=-x +
%  4$ is the point-slope equation.
%\end{badexample}
  

\begin{goodexample}
  
We first check that $(1,3)$ is a point on the curve by plugging these
values in: $3=1+2/1$. 
  The derivative of the curve $y = x + \frac{2}{x}$ is
%  \begin{displaymath}
\begin{equation}
    \frac{dy}{dx} = 1 - \frac{2}{x^2}.
%\label{eqn:refexample}
\end{equation}
%  \end{displaymath}
  A line tangent to the curve at the point $(1,3)$ will have slope
  \begin{displaymath}
    \left. \frac{dy}{dx} \right| _{x=1} = 1 - \frac{2}{(-1)^2} = -1.
  \end{displaymath}
  Using the point-slope formula with $m = -1$, $x_0 = 1$, and $y_0 = 3$ gives
  the formula for the line $y - 3 = -1(x-1)$.  Solving for y and
  simplifying gives
  \begin{displaymath}
    y = -x + 4.
  \end{displaymath}
  This is the equation of the line tangent to the curve at that point.
\end{goodexample}  


%%%%%%%%%%%%%%%%%%%%%%%%%%%%%%%%%%%%%%%%%%%%%%%%%%%%%%%%%%%%%%%%%%%%%%%%%
%\subsubsection*{Vague and uninformative sentences}
    
\item
\begin{badexample}
      We use calculus to find that $y = 3x^2+1$ has a slope of $3$ at
      $x = 1/2$.
\end{badexample}

\begin{goodexample}
    To find the slope of the curve $y = 3x^2+1$ at the point $x = 1/2$
    we find $\frac{dy}{dx}$ evaluated at $x = 1/2$. We compute
    \begin{eqnarray*}
      \frac{dy}{dx} &  = &  6x,  \\  
      \Rightarrow \left. \frac{dy}{dx} \right| _{x=1/2} & = & 6 \left( \frac{1}{2}
      \right) = 3. 
    \end{eqnarray*}
\end{goodexample}

\item
                 %not sure if this is a good example of
                 %not, seems a little contrived 
\begin{badexample}
    \begin{eqnarray*}
      \frac{d}{dx}\left( (x^2 + 1) (x^3 +3) \right) &  = & 
      (x^2+1)3x^2 +(x^3+3)2x \\ 
      &  = & 5x^4+3x^2+6x
    \end{eqnarray*}
    is the derivative.
\end{badexample}
    
\begin{goodexample}
    To take the derivative of a product of $2$ functions, we use the
    product rule, $(fg)'=f'g+fg'$.
In our case we have
    \begin{eqnarray*}
      \frac{d}{dx}\left( (x^2 + 1) (x^3 +3) \right) &  = & 
      \left(\frac{d}{dx} (x^2 + 1) \right) (x^3 +3)
	+(x^2 + 1) \left(\frac{d}{dx} (x^3 +3)\right)\\
	& =& (2x)(x^3+3) +(x^2 + 1)(3x^2)\\
      &  = & 5x^4+3x^2+6x.
    \end{eqnarray*}
\end{goodexample}

%%%%%%%%%%%%%%%%%%%%%%%%%%%%%%%%%%%%%%%%%%%%%%%%%%%%%%%%%%%%%%%%%%%%%%%%%
%\subsubsection*{Referring to previous equations and figures.}
\item
You may have noticed that one of the equations above has been labeled
equation number one by putting ``(1)'' at the right hand
margin. If you need to refer back to an equation or figure, label it
and then refer to it by its label. Do not draw arrows.

\begin{badexample}
Using the equation from before, the slope is $-1$. \fbox{\rm Which equation?}
\end{badexample}

\begin{goodexample}
Using (1), the derivative at $x=1$ is $-1$.
\end{goodexample}
\end{itemize}

%\marginpar{conclusion.}






